% Generated draft: requires mathematical and source review.
\documentclass[11pt]{article}
\usepackage[margin=25mm]{geometry}
\usepackage{fontspec}
\setmainfont{Latin Modern Roman}
\usepackage{amsmath,amssymb,mathtools,unicode-math}
\setmathfont{Latin Modern Math}
\usepackage{enumitem,needspace,xurl,hyperref}
\setlength{\parindent}{0pt}
\setlength{\parskip}{5pt}
\setlength{\emergencystretch}{3em}
\newcommand{\N}{\mathbb{N}}
\newcommand{\Z}{\mathbb{Z}}
\newcommand{\Q}{\mathbb{Q}}
\newcommand{\R}{\mathbb{R}}
\newcommand{\C}{\mathbb{C}}
\newcommand{\F}{\mathbb{F}}
\newcommand{\A}{\mathbb{A}}
\newcommand{\PP}{\mathbb{P}}
\newcommand{\E}{\mathbb{E}}
\newcommand{\eps}{\varepsilon}
\newcommand{\dd}{\,\mathrm{d}}
\begin{document}
\section*{UnsolvedMath ID: 30002255}
\textbf{Source problem code:} catalogue entry
\subsection{Definitions and mathematical statement}
In the usual monster-model setting, let \(T\) be a complete first-order theory and let \(\mathfrak C\models T\) be a sufficiently saturated and strongly homogeneous model. For tuples of a fixed small length, ordinary Lascar equivalence is defined by
\[
 a\equiv_L b
 \quad\Longleftrightarrow\quad
 (\exists\sigma\in\operatorname{Autf}(\mathfrak C))\;\sigma(a)=b,
\]
where \(\operatorname{Autf}(\mathfrak C)\) is the subgroup generated by the pointwise stabilizers of small elementary submodels of \(\mathfrak C\). This defines ordinary Lascar equivalence; the supplied record does not define its notation \(\equiv_L^\alpha\) or explain the meaning of the level \(\alpha\).

For equivalence relations \(E\) on a standard Borel space \(X\) and \(F\) on a standard Borel space \(Y\), write \(E\leq_B F\) if there is a Borel map \(f:X\to Y\) such that
\[
 (\forall x,x'\in X)\qquad
 x\mathrel{E}x'\quad\Longleftrightarrow\quad f(x)\mathrel{F}f(x').
\]
The relations have the same Borel cardinality if \(E\leq_B F\) and \(F\leq_B E\); write \([E]_{\sim_B}\) for this equivalence class. This comparison requires specified standard Borel domains, but does not itself require \(E\) and \(F\) to be Borel subsets of their product spaces.

\textbf{Problem as supplied.} For the level-\(\alpha\) Lascar equivalence relation \(\equiv_L^\alpha\) in a complete first-order theory, determine which Borel cardinalities \([\equiv_L^\alpha]_{\sim_B}\) can occur.

\textbf{Unresolved qualifications.} The source record does not specify the definition or admissible values of \(\alpha\), the underlying tuple domain, the standard Borel coding, countability assumptions on the language or parameters, or which of these data and \(T\) are fixed or allowed to vary. It also does not specify whether the relation is considered globally or restricted to a type or a Kim--Pillay strong type. These omissions prevent a uniquely quantified classification problem; the definition of ordinary Lascar equivalence above does not resolve them.
\subsection{Short English statement}
Classify the possible Borel cardinalities of the source's level-\(\alpha\) Lascar equivalence relations. The meaning of the level, the domains and coding, and the parameters allowed to vary require clarification. The supplied record labels the problem partially solved; this is a statement-only draft requiring human mathematical review.
\subsection{Sources}
\begin{itemize}
\item Local UnsolvedMath registry, \texttt{data/unsolvedmath/problems.json}, supplied record with permanent ID \(30002255\), source code \texttt{catalogue entry}, titled \emph{Borel Cardinalities of Lascar Equivalence Relations}. The record gives proposed year 2013 and names no proposer.
\item \emph{Model Theory: Groups, Geometry, and Combinatorics} (2013), Oberwolfach Reports, pp.~5--66, DOI 10.4171/owr/2013/01. \url{https://doi.org/10.4171/owr/2013/01}. The supplied record describes the report as documenting workshop discussions and says that this question is its fifth explicit descriptive-set-theoretic problem. Its statement-status label is \texttt{corrected\_verified}; that label and verification claim are supplied curation, not independent verification here.
\item \emph{The Borel Cardinality of Lascar Strong Types}, attributed in the supplied assessment to Kaplan, Miller, and Simon. \url{https://doi.org/10.1112/jlms/jdu041}.
\item Supplied literature assessment dated \texttt{2026-08-21}: status \texttt{partially\_solved}. It attributes to Kaplan, Miller, and Simon non-smoothness whenever the Lascar relation on a Kim--Pillay strong type is nontrivial, while stating that this does not classify all possible Borel cardinalities. Here smoothness means Borel reducibility to equality on \(\R\). The assessment lists both references above as checked. This is supplied dated curation derived from the Oberwolfach material, not an independently verified literature assessment; its precise theorem hypotheses and time-sensitive status require independent review.
\end{itemize}
\end{document}
